\section{Exactness of the real relaxation for two-zero weighing completion}
\label{sec:real-weighing-exactness}

The eight-state completion test admits a stronger interpretation.
For a weighing matrix with two zeros per row at an admissible order,
allowing arbitrary real entries in the holes does not enlarge the
class of matrices that can be completed. A real completion either
already supplies a sign state or produces a conference matrix from
which a sign completion can be recovered. We also classify the
continuous solution families and sharpen the eight-fill bound.

Throughout $4\mid n$, $n\ge4$, and $W$ is $W(n,n-2)$. Use the
previous pass's row and column normalization: all zero positions
are diagonal two-square blocks, and
\[
 F=\begin{pmatrix}1&1\\1&-1\end{pmatrix},\qquad
 A_{ij}=\tfrac12W_{ij}F\in\mathcal D\quad(i\ne j),
\]
where $\mathcal D$ is the eight-element signed-permutation group.
Put
\[
 I=\begin{pmatrix}1&0\\0&1\end{pmatrix},\quad
 J=\begin{pmatrix}0&1\\-1&0\end{pmatrix},\quad
 D=\begin{pmatrix}1&0\\0&-1\end{pmatrix},\quad
 S=\begin{pmatrix}0&1\\1&0\end{pmatrix}.
\]
Let $\mathcal R=\operatorname{span}_{\mathbb R}\{I,J\}$ and
$\mathcal F=\operatorname{span}_{\mathbb R}\{D,S\}$.
The orthogonal group $O(2)$ consists exactly of the unit coefficient
circles in these two planes.

\subsection{A linear kernel controls every real completion}
A real zero-supported completion has diagonal fill blocks
$X_i=R_iF$, $R_i\in O(2)$, and satisfies
\[
 A_{ij}R_j^{\mathsf T}+R_iA_{ji}^{\mathsf T}=0.
\]
For a root $R=R_0$, define
\[
 R_j=-A_{j0}R^{\mathsf T}A_{0j}\quad(j>0).
\]
Let $K\subseteq M_2(\mathbb R)$ be the homogeneous linear space
where all remaining cross equations hold. A root gives a real
completion exactly when $R\in K\cap O(2)$; it gives a Hadamard
filling exactly when $R\in K\cap\mathcal D$.

\begin{lemma}[The only possible component sections]
Before imposing any additional consequences of the weighing
property, each of $K\cap\mathcal R$ and $K\cap\mathcal F$ is zero,
the whole plane, a coordinate-axis line, or one of the two diagonal
lines in the displayed coefficient bases.
\end{lemma}
\begin{proof}
After multiplying the cross equation for $i,j>0$ on the left and
right by signed permutations, it becomes
\[
 R+L_{ij}R^{\mathsf T}Q_{ij}=0,
\]
where
\[
 L_{ij}=A_{0j}A_{ij}^{\mathsf T}A_{i0},\qquad
 Q_{ij}=A_{0i}A_{ji}^{\mathsf T}A_{j0}
\]
are signed permutations. The operator
$R\mapsto-L_{ij}R^{\mathsf T}Q_{ij}$ is Frobenius-orthogonal.
It either interchanges the two planes, forcing their pure-component
solutions to zero, or preserves them and acts on each coefficient
pair by a signed permutation. The fixed space of a two-dimensional
signed permutation is zero, a full plane, an axis line, or a
diagonal line. Intersecting these fixed spaces introduces no other
line directions. This proves the assertion.
\end{proof}

\subsection{A diagonal real state produces a conference completion}
The next argument invokes the classical conference normalization
theorem, not an unproved rounding principle. Geramita--Verner,
\emph{Orthogonal Designs with Zero Diagonal}, Theorem 2,
\url{https://doi.org/10.4153/CJM-1976-028-3}, restates the
Delsarte--Goethals--Seidel result: a conference matrix at order
$0\bmod4$ can be made skew by sign changes. Equivalently its
unique zero per row and column can be filled by a signed permutation
$Y$ so that $U+Y$ is Hadamard. Both fills $Y$ and $-Y$ work.

\begin{lemma}[Conference conversion]
If $K$ contains a diagonal line in either coefficient plane, then
that entire plane is contained in $K$.
\end{lemma}
\begin{proof}
Normalize a diagonal root to an orthogonal matrix $R$. Its entries
are all $\pm1/\sqrt2$. Every propagated $R_i$ has the same entry
form, so $R_iF$ has exactly one nonzero $\pm\sqrt2$ in each row
and column. Thus the real fill is $X=\sqrt2T$, where $T$ is a
block diagonal signed permutation supported on half the zero
positions of $W$.

The completion equations give
$WT^{\mathsf T}+TW^{\mathsf T}=0$. Hence
$C=WT^{\mathsf T}$ is skew, has diagonal two-square blocks zero,
and satisfies $CC^{\mathsf T}=(n-2)I$. The integer matrix
\[
 U=C+I_n
\]
has entries $0,\pm1$ and $UU^{\mathsf T}=(n-1)I$: it is a
conference matrix, with its single remaining zero in each row at
the off-diagonal position of the corresponding coordinate pair.
Apply the classical conference theorem, first permuting its zero
matching onto the diagonal if needed. It supplies a signed
permutation $Y$ on these remaining zeros such that both $U\pm Y$
are Hadamard. The paired diagonal blocks of $U$ are $I_2$, so
orthogonality forces $Y_i^{\mathsf T}=-Y_i$, that is $Y_i=\pm J$.

Undoing $T$ gives the two sign completions
\[
 H_\pm=(U\pm Y)T=W+(I_n\pm Y)T.
\]
The matrices $T$ and $YT$ have complementary supports inside each
original zero block. Thus each hole receives one sign and every
nonzero entry of $W$ remains unchanged. At the root the two fill
states are
\[
 R_\pm=\frac{(I_2\pm Y_0)T_0F}{2}.
\]
They lie in the same orientation plane as the original diagonal
root $T_0F/\sqrt2$, and are linearly independent because $Y_0=\pm J$
has no real eigenvector. Both belong to $K$. Their span is therefore
the entire component plane, as asserted.
\end{proof}

\begin{theorem}[Exact real feasibility and solution topology]
A real matrix $H$ with $HH^{\mathsf T}=nI_n$, obtained by changing
only zeros of $W$, exists if and only if a zero-preserving Hadamard
filling exists. More precisely each component section
$K\cap\mathcal R$, $K\cap\mathcal F$ is only zero, a coordinate
axis, or the full plane. Each axis contributes two isolated
completions, both sign-valued. Each full plane contributes a full
circle of real completions containing exactly four sign fillings.
There are no isolated nonsign real completions.
\end{theorem}
\begin{proof}
The preceding lemmas exclude an isolated diagonal-line section.
Every nonzero permitted section intersects $O(2)$ and contains
signed-permutation roots. Conversely a real orthogonal root lies
in one of these sections. The counts follow by intersecting an
axis or a plane with its unit coefficient circle, and the root
propagation is injective. A nonsign real root can therefore occur
only on one of the full circles, not as an isolated solution.
\end{proof}

This is an exactness theorem for this specific real feasibility
relaxation. It does not claim that arbitrary approximate real
fillings can be stably rounded, or that real feasibility always
holds. The unfillable order-twelve example has both sections zero.

\subsection{At most six fillings beyond order four}
\begin{proposition}[A sharper cardinality bound]
For $n\ge8$ the number of zero-preserving Hadamard fillings belongs
to $\{0,2,4,6\}$ and is at most six. The bound six is attained.
At order four there are exactly eight fillings.
\end{proposition}
\begin{proof}
Each component contributes zero, two or four sign roots. If both
planes were full, $K$ would be all of $M_2(\mathbb R)$. At $n\ge8$
there is a block triangle and hence at least one equation
$R+LR^{\mathsf T}Q=0$. It cannot hold identically. Substituting
$R=I$ would give $Q=-L^{\mathsf T}$ and thus
$R=LR^{\mathsf T}L^{\mathsf T}$ for every matrix. This map reverses
products; if it were the identity, all two-square matrices would
commute, a contradiction. Thus $K$ is proper and eight is impossible.
The skew order-four conference tensor with $F$ has one full plane
and one axis, attaining six at order eight. For $n=4$ there are
only two blocks and no triangle constraints; every one of the
eight roots works. The displayed set gives necessary possibilities;
no assertion that every listed value occurs at every order is made.
\end{proof}

\subsection{Exact checks and a nonsign circle point}
The verifier checks the real kernel sections for the four basic
cases: order four has two full planes and eight fillings; the
skew-conference tensor at order eight has an axis and a full plane,
giving six; the symmetric-conference tensor at order twelve has
one axis and two fillings; the inconsistent order-twelve example
has no real or sign completion. It also exhausts all 256 normalized
$W(8,6)$ matrices with standard zero support and first-row/sign
normalizations. Every case has one axis and one full plane and
exactly six sign fillings.

For a concrete nonsign real completion, the order-eight full
reflection plane contains the rational root $(3D+4S)/5$.
Its resulting completion has exact Gram $8I$ and nonsign hole
entries. The converter is separately tested at diagonal states:
it constructs the two conference fills and checks both resulting
Hadamard matrices and the independence of their root states using
only integer and rational arithmetic. No isolated diagonal-line
example was found; the conversion theorem explains why none can
occur. No novelty or all-order Hadamard existence claim is made.
